% Density and full cluster tree, same curve/topology as the defense.
\begin{tikzpicture}[every node/.style={font=\fontsize{10}{11}\selectfont},line cap=round]
\useasboundingbox (-0.7,-0.65) rectangle (13.8,3.9);
\def\curve{(0.2,0.2)..controls(0.7,0.2)and(0.9,2)..(1.2,2)
..controls(1.5,2)and(1.6,0.9)..(1.8,0.9)
..controls(2,0.9)and(2.2,1.5)..(2.4,1.5)
..controls(2.5,1.5)and(2.6,1.3)..(2.7,1.3)
..controls(2.9,1.3)and(3.1,3.2)..(3.4,3.2)
..controls(3.7,3.2)and(3.9,0.6)..(4.2,0.6)
..controls(4.4,0.6)and(4.6,1.3)..(4.8,1.3)
..controls(5,1.3)and(5.4,0.2)..(5.8,0)}
\begin{scope}
\draw[-{Latex[length=2mm]},inriagrayblue] (0,0)--(6,0) node[right] {$y$};
\draw[-{Latex[length=2mm]},inriagrayblue] (0,0)--(0,3.5) node[above] {$f(y)$};
\begin{scope}\clip (0,1.1) rectangle (6,3.4);
\fill[inriablue!20] \curve -- (5.8,0)--(.2,0)--cycle;\end{scope}
\draw[inriablue,line width=1.3pt] \curve;
\draw[dashed,inriagrayblue] (0,1.1)--(6,1.1) node[right] {$\lambda$};
\foreach \a/\b in {.76/1.60,1.98/4.00,4.52/5.00}
  \draw[inriablue,line width=2.5pt] (\a,-.12)--(\b,-.12);
\node[text=inriagrayblue] at (2.9,-.55) {Density superlevel sets};
\end{scope}
\begin{scope}[xshift=7.5cm]
\draw[-{Latex[length=2mm]},inriagrayblue] (0,0)--(0,3.5) node[above] {$\lambda$};
\draw[inriared,line width=1.35pt] (1.2,2)--(1.2,.9)--(2.05,.9)--(2.05,.6)--(3.42,.6)--(3.42,.22);
\draw[inriared,line width=1.35pt] (2.4,1.5)--(2.4,1.3)--(2.9,1.3)--(2.9,.9)--(2.05,.9);
\draw[inriared,line width=1.35pt] (3.4,3.2)--(3.4,1.3)--(2.9,1.3);
\draw[inriared,line width=1.35pt] (4.8,1.3)--(4.8,.6)--(3.42,.6);
\draw[dashed,inriagrayblue] (0,1.1)--(5.4,1.1);
\foreach \x in {1.2,2.9,4.8}\fill[inriablue] (\x,1.1) circle[radius=.07];
\node[text=inriagrayblue] at (3.05,-.55) {The same clusters in the tree};
\end{scope}
\end{tikzpicture}
